%% problem_id: S014-lurie-ultracategories-prop-3.4.6
%% source_id: S014 (Jacob Lurie, Ultracategories)
%% source_file: lurie-conceptual.pdf
%% statement_locator: p. 36, Proposition 3.4.6
%% proof_locator: pp. 38--41
%% representation: PDF; transcribed from rendered page images (pdftoppm -r 150 -png), not from the text layer
%% status: TRANSCRIPTION ONLY (difficulty judgement is made separately, by the Lead)
%%
%% Transcription notes (every symbol-level reading below was checked against magnified crops of the page
%% images; nothing was taken from the PDF text layer):
%%  * Pages transcribed, in page order: 36, 38, 39, 40, 41. The mathematical content of each of those
%%    pages is reproduced verbatim, including the items standing around Proposition 3.4.6 (Theorem 3.4.4,
%%    Variant 3.4.5, Corollaries 3.4.7--3.4.9) and the lemmas of section 3.5, since the proof of
%%    Proposition 3.4.6 refers to them and they lie on the supplied pages.
%%  * Boxed cross-reference numbers of the source (hyperref boxes) are transcribed as plain references.
%%  * Places where the printed text deviates from what the mathematics would suggest are kept exactly as
%%    printed and flagged by a %% [NOTE: ...] line; nothing has been silently corrected.
%%  * No passage of the supplied pages was illegible, so there are no [ILLEGIBLE: ...] markers.

\documentclass[11pt]{article}
\usepackage{amsmath}
\usepackage{amssymb}
\usepackage{amsthm}
\usepackage{mathrsfs}
\usepackage[margin=1in]{geometry}

\theoremstyle{plain}
\newtheorem{theorem}{Theorem}
\newtheorem{proposition}{Proposition}
\newtheorem{lemma}{Lemma}
\newtheorem{corollary}{Corollary}

\theoremstyle{definition}
\newtheorem{variant}{Variant}
\newtheorem{notation}{Notation}

\begin{document}

%% --- page 36 ---

%% [NOTE: what follows is the tail of a proof whose \begin{proof} stands on printed page 35 (not among the
%%  supplied images); it is reproduced here because it occupies the top of printed page 36.]
\begin{proof}[Proof (continued from the preceding page)]
Let $\mathscr{G}$ denote the sheaf of sets on $X$ given by the formula $\mathscr{G}(U) = \prod_{x \in U} G(x)$. By construction, the presheaf $\mathscr{F}_G$ is contained in $\mathscr{G}$. It will therefore suffice to show that if $\{U_\alpha\}$ is a collection of open subsets of $X$ having union $U = \bigcup_\alpha U_\alpha$, and $\phi_\bullet \in \mathscr{G}(U)$ is a section having the property that its image in each $\mathscr{G}(U_\alpha)$ belongs to $\mathscr{F}_G(U_\alpha)$, then $\phi_\bullet$ belongs to $\mathscr{F}_G(U)$. Write $\phi_\bullet = \{\phi_x\}_{x \in U}$ for some collection of elements $\phi_x \in \mathscr{F}_x$. Suppose that $f : S \to U \subseteq X$ is a map of sets and that $\mu$ is an ultrafilter on $S$ with the property that $\int_S f(s)d\mu$ belongs to $U$; we wish to prove the identity
\[
q_\mu(\{\phi_{f(s)}\}_{s \in S}) = \sigma_\mu(\phi_{\int_S f(s)d\mu})
\]
in the ultraproduct $\int_S G(f(s))d\mu$. Choose an index $\alpha$ such that $\int_S f(s)d\mu$ belongs to the subset $U_\alpha \subseteq U$. Set $S_0 = \{s \in S : f(s) \in U_\alpha\}$. Then $\mu(S_0) = 1$, so we can replace $S$ by $S_0$ and $f$ by $f|_{S_0}$. In this case, the relevant identity follows from our assumption that $\{\phi_x\}_{x \in U_\alpha}$ belongs to $\mathscr{F}_G(U_\alpha)$.
\end{proof}

We can now give a precise statement of our main result.

\renewcommand{\thetheorem}{3.4.4}
\begin{theorem}
Let $X$ be a compact Hausdorff space. Then the construction $G \mapsto \mathscr{F}_G$ induces an equivalence of categories $\mathrm{Fun}^{\mathrm{LUlt}}(X,\mathrm{Set}) \to \mathrm{Shv}(X)$.
\end{theorem}

\renewcommand{\thevariant}{3.4.5}
\begin{variant}
Recall that a category $\mathcal{M}$ is said to be \emph{compactly generated} if it admits small colimits and is generated under filtered colimits by a small collection of compact objects. Equivalently, $\mathcal{M}$ is compactly generated if there exists an equivalence $\mathcal{M} \simeq \mathrm{Fun}^{\mathrm{lex}}(\mathcal{C},\mathrm{Set})$, for some small category $\mathcal{C}$ which admits finite limits (which can then be identified with the opposite of the category of compact objects of $\mathcal{M}$). If $\mathcal{M}$ is compactly generated and $X$ is a compact Hausdorff space, then Theorem 3.4.4 determines equivalence from the category of left ultrafunctors $\mathrm{Fun}^{\mathrm{LUlt}}(X,\mathcal{M})$ to the category $\mathrm{Shv}(X;\mathcal{M})$ of $\mathcal{M}$-valued sheaves on $X$, given by the composition
\begin{align*}
\mathrm{Fun}^{\mathrm{LUlt}}(X,\mathcal{M}) \;&\simeq\; \mathrm{Fun}^{\mathrm{LUlt}}(X,\mathrm{Fun}^{\mathrm{lex}}(\mathcal{C},\mathrm{Set}))\\
&\simeq\; \mathrm{Fun}^{\mathrm{lex}}(\mathcal{C},\mathrm{Fun}^{\mathrm{LUlt}}(X,\mathrm{Set}))\\
&\to\; \mathrm{Fun}^{\mathrm{lex}}(\mathcal{C},\mathrm{Shv}(X))\\
&\simeq\; \mathrm{Shv}(X;\mathrm{Fun}^{\mathrm{lex}}(\mathcal{C},\mathrm{Set}))\\
&\simeq\; \mathrm{Shv}(X;\mathcal{M}).
\end{align*}
\end{variant}

One can also describe this equivalence directly, using a variant of Construction 3.4.2.

Let $X$ be a compact Hausdorff space containing a point $y$ and let $G : X \to \mathrm{Set}$ be a left ultrafunctor. For every open subset $U \subseteq X$ containing $y$, the construction $\{\phi_x\}_{x \in U} \mapsto \phi_y$ determines a map $\mathscr{F}_G(U) \to G(y)$. These maps are compatible as $U$ varies, and therefore induce a map of sets $\mathscr{F}_{G,y} \to G(y)$. The main ingredient in our proof of Theorem 3.4.4 is the following result, whose proof we defer to §3.5.

\renewcommand{\theproposition}{3.4.6}
\begin{proposition}
Let $X$ be a compact Hausdorff space and let $G : X \to \mathrm{Set}$ be a left ultrafunctor. Then, for each point $y \in X$, the preceding construction induces a bijection $\mathscr{F}_{G,y} \to G(y)$.
\end{proposition}

\renewcommand{\thecorollary}{3.4.7}
\begin{corollary}
Let $X$ be a compact Hausdorff space and let $\alpha : G \to H$ be a natural transformation of left ultrafunctors $G, H : X \to \mathrm{Set}$. Then $\alpha$ is an isomorphism if and only if, for each open subset $U \subseteq X$, the induced map $\mathscr{F}_G(U) \to \mathscr{F}_H(U)$ is an isomorphism.
\end{corollary}

\renewcommand{\thecorollary}{3.4.8}
\begin{corollary}
Let $X$ be a compact Hausdorff space. Then the functor $\mathrm{Fun}^{\mathrm{LUlt}}(X,\mathrm{Set}) \to \mathrm{Shv}(X)$ preserves small colimits.
\end{corollary}

\begin{proof}
Since small colimits $\mathrm{Shv}(X)$ are computed stalkwise, it will suffice to show that for each point $x \in X$, the functor $G \mapsto \mathscr{F}_{G,x}$ commutes with small colimits. By virtue of Proposition 3.4.6, this is equivalent to the evaluation functor $G \mapsto G(x)$, which preserves small colimits by virtue of Remark 1.4.3.
\end{proof}

\renewcommand{\thecorollary}{3.4.9}
\begin{corollary}
Let $X$ be a compact Hausdorff space and let $G : X \to \mathrm{Set}$ be a left ultrafunctor. Then every subsheaf of $\mathscr{F}_G$ has the form $\mathscr{F}_{G_0}$, for some uniquely determined subobject $G_0 \subseteq G$ in the category of left ultrafunctors $\mathrm{Fun}^{\mathrm{LUlt}}(X,\mathrm{Set})$.
\end{corollary}

%% --- page 38 ---

%% [NOTE: the paragraph beginning "We now show ..." belongs to the argument of section 3.4 (the equivalence
%%  (a) <=> (b) discussed there) and not to section 3.5; section 3.5, "The Proof of Proposition 3.4.6",
%%  begins lower on this same page. Both are transcribed because both stand on the supplied page.]
We now show that the sheaf $\mathscr{F}_G$ is constant in some neighborhood of any point $y \in X$. Let $n$ be the cardinality of the stalk $\mathscr{F}_{G,y}$. Choose an open neighborhood $U$ of the point $y$ and a finite collection of sections $\phi_{1,\bullet},\dots,\phi_{n,\bullet} \in \mathscr{F}_G(U)$ having distinct images in $\mathscr{F}_{G,y}$. Let us identify each $\phi_{i,\bullet}$ with a tuple $\{\phi_{i,x} \in G(x)\}_{x \in U}$ satisfying condition $(\ast)$ of Construction 3.4.2. Note that, since $G$ is an ultrafunctor, condition $(\ast)$ has the following consequence:

\smallskip
\noindent $(\ast')$ Let $f : S \to U$ be a map of sets and let $\mu$ be an ultrafilter on $S$ such that $\int_S f(s)d\mu$ belongs to $\mu$. If $\phi_{i,f(s)} = \phi_{j,f(s)}$ for all $s \in S$, then $\phi_{i,\int_S f(s)d\mu} = \phi_{j,\int_S f(s)d\mu}$.
\smallskip

%% [NOTE: in the printed line of $(\ast')$ above the symbol after "belongs to" is a lowercase mu (not a
%%  capital U); it is transcribed as printed.]
Let $V \subseteq U$ be an open neighborhood of $y$ whose closure $\overline{V}$ is contained in $U$. For $1 \le i < j \le n$, let $K_{i,j} \subseteq \overline{V}$ denote the subset consisting of those points $x$ for which $\phi_{i,x} = \phi_{j,x}$. It follows from $(\ast')$ that each of the sets $K_{i,j}$ is closed. By construction, the sets $K_{i,j}$ do not contain the point $y$.

Let $K' \subseteq \overline{V}$ be the set of all points $x$ for which the set $G(x)$ contains some element $\psi_x \notin \{\phi_{1,x},\dots,\phi_{n,x}\}$. Note that, for every ultrafilter $\mu$ on the set $K'$, the image $\{\psi_x\}_{x \in K'}$ under the map
\[
\prod_{x \in K'} G(x) \to \int_{K'} G(x)d\mu \xrightarrow{\ \sigma_\mu^{-1}\ } G(\int_{K'} x d\mu)
\]
does not belong to the set $\{\phi_{1,\int_{K'} x d\mu},\dots,\phi_{n,\int_{K'} x d\mu}\}$. It follows that $K'$ is a closed subset of $\overline{V}$, which also does not contain the point $y$.

Let $W \subseteq V$ be an open neighborhood of $y$ which is disjoint from the closed sets $K'$ and $K_{i,j}$. For each point $x \in W$, the elements $\phi_{1,x},\dots,\phi_{n,x}$ are pairwise distinct and exhaust $G(x)$. It follows from Proposition 3.4.6 that the germs of the sections $\phi_1,\dots,\phi_n$ are pairwise distinct and exhaust the stalk $\mathscr{F}_{G,x}$. Consequently, they determine an isomorphism of $\mathscr{F}_G|_W$ is isomorphic with the constant sheaf associated to the finite set $\{1,2,\dots,n\}$. This completes the proof that $(a)$ implies $(b)$.

Now suppose that $\mathscr{F}_G$ is locally constant; we wish to show that $G$ is an ultrafunctor. Choose a map $f : S \to X$ and an ultrafilter $\mu$ on $S$, and set $x = \int_S f(s)d\mu$. We must show that the left ultrastructure map $\sigma_\mu : G(x) = G(\int_S f(s)d\mu) \to \int_S G(f(s))d\mu$ is bijective. Choose an open neighborhood $U$ of $x$ such that $\mathscr{F}_G|_U$ is isomorphic to the constant sheaf $\underline{J}_U$, for some finite set $J$. Let $V$ be an open neighborhood of $x$ whose closure is contained in $U$, and set $S_0 = f^{-1}(V)$. Since $x$ belongs to $V$, we must have $\mu(S_0) = 1$. We may therefore replace $S$ by $S_0$ and $\mu$ by its restriction to $S_0$, and thereby reduce to the case where the function $f$ takes values in $V$. Replacing $X$ by the compact set $\overline{V} \subseteq X$, we can reduce to the case where $\mathscr{F}_G = \underline{J}_X$ is itself the constant sheaf associated to a finite set $J$. In this case, we have an isomorphism of sheaves $\mathscr{F}_G \simeq \mathscr{F}_{G'}$ where $G' : X \to \mathrm{Set}$ is the constant ultrafunctor taking the value $J$. It follows from Theorem 3.4.4 that that $G$ is isomorphic to $G'$ (as a left ultrafunctor), and is therefore also an ultrafunctor. \qed

%% [NOTE: the printed sentence reads "that that $G$ is isomorphic to $G'$"; the doubling stands in the source.]
\smallskip

\noindent\textbf{3.5. The Proof of Proposition 3.4.6.} Throughout this section, we fix a compact Hausdorff space $X$ and a left ultrafunctor $G : X \to \mathrm{Set}$. Our goal is to compute the stalks of the sheaf $\mathscr{F}_G$ of Construction 3.4.2. We begin by establishing a weak version of Proposition 3.4.6.

\renewcommand{\thelemma}{3.5.1}
\begin{lemma}
For each point $y \in X$, the map $\mathscr{F}_{G,y} = \varinjlim_{y \in U} \mathscr{F}_G(U) \to G(y)$ of Proposition 3.4.6 is injective.
\end{lemma}

\begin{proof}
Let $U$ be an open neighborhood of the point $y$ and suppose we are given a pair of elements $\phi_\bullet = \{\phi_x\}_{x \in U}$ and $\psi_\bullet = \{\psi_x\}_{x \in U}$ of $\mathscr{F}_G(U)$ satisfying $\phi_y = \psi_y$ (as elements of the set $G(y)$). We wish to show that there exists an open subset $V \subseteq U$ containing the point $y$ such that $\phi_x = \psi_x$ for each $x \in V$. Assume otherwise. Set $U_0 = \{x \in U : \phi_x \neq \psi_x\} \subseteq U$, and let $\mathcal{U}$ be the collection of all subsets of $U$ having the form $U_0 \cap V$, where $V$ is an open neighborhood of $y$. Then $\mathcal{U}$ is closed under finite intersections and does not contain the empty set. Applying Proposition 1.1.10 we deduce that there exists an ultrafilter $\mu$ on the set $U$ such that $\mu(U_0 \cap V) = 1$, whenever $V$ is an open neighborhood of $y$. We then have $\int_U x d\mu = y$. Invoking our assumption that $\phi_y = \psi_y$ and the definition of the sheaf $\mathscr{F}_G$, we deduce that $\phi_\bullet$ and $\psi_\bullet$ have the same image in the ultraproduct $\int_U G(x)d\mu$. It follows that $\mu(U_0) = 0$, contradicting the definition of $U_0$.
\end{proof}

%% --- page 39 ---

To complete the proof of Proposition 3.4.6 we must show that each element $\varphi_y \in G(y)$ can be extended to an element $\phi_\bullet \in \mathscr{F}_G(U)$ for some open neighborhood $U$ of the point $y$. The proof will require some auxiliary constructions.

\renewcommand{\thenotation}{3.5.2}
\begin{notation}
Choose a set $T$ equipped with a bijection $f_0 : T \to X$. We regard $f_0$ as a continuous bijection of topological spaces, where $T$ is endowed with the discrete topology. We will use the letter $\nu$ to denote a typical ultrafilter on $T$: that is, a point in the Stone-Čech compactification $\beta T$. Let $i : T \hookrightarrow \beta T$ be the canonical embedding, which assigns to each point $t \in T$ the corresponding principal ultrafilter $i(t) = \delta_t$. By virtue of Proposition 3.2.7, the map $f_0 : T \to X$ extends uniquely to a continuous map $f : \beta T \to X$, given concretely by the formula $f(\nu) = \int_T f_0(t)d\nu$.

Let $\mathscr{G}_0$ denote the sheaf of sets on the (discrete) topological space $T$ whose stalks are given by $\mathscr{G}_{0t} = G(f_0(t))$, and let $\mathscr{G}$ denote the direct image $i_*\mathscr{G}_0$. Then $\mathscr{G}$ is a sheaf of sets on $\beta T$ whose value on closed and open sets is given by the formula
\[
\mathscr{G}(\beta T_0) = \prod_{t \in T_0} G(f_0(t)),
\]
and whose stalk at a point $\nu \in \beta T$ is given by the formula $\mathscr{G}_\nu \simeq \int_T G(f_0(t))d\nu$.

Choose a set $S$ and a bijection $g_0 : S \to \beta T$, which we identify with a collection of ultrafilters $\{\nu_s\}_{s \in S}$ on the set $T$. We regard $g_0$ as a continuous bijection of topological spaces, where $S$ is equipped with the discrete topology and $\beta T$ with the topology of Construction 3.2.2. We will use the letter $\mu$ to denote a typical ultrafilter on $S$: that is, a point in the Stone-Čech compactification $\beta S$. Let $j : S \hookrightarrow \beta S$ be the canonical embedding, which assigns to each point $s \in S$ the corresponding principal ultrafilter $j(s) = \delta_s$. Let $\mathscr{H}_0 = g_0^*\mathscr{G}$ denote the sheaf of sets on $S$ whose stalks are given by $\mathscr{H}_{0s} = \int_T G(f_0(t))d\nu_s$. Let $\mathscr{H}$ denote the sheaf of sets on $\beta S$ given by the direct image $j_*\mathscr{H}_0$. Then $\mathscr{H}$ is given on closed and open sets by the formula
\[
\mathscr{H}(\beta S_0) = \prod_{s \in S_0} \int_T G(f_0(t))d\nu_s,
\]
and its stalk at a point $\mu \in \beta S$ is given by $\mathscr{H}_\mu \simeq \int_S (\int_T G(f_0(t))d\nu_s)d\mu$.

By virtue of Proposition 3.2.7 the bijection $g_0 : S \to \beta T$ admits a unique continuous extension $g : \beta S \to \beta T$, given concretely by the formula $g(\mu) = \int_S \nu_s d\mu$ (see Example 3.2.8). We have a canonical isomorphism $j^*g^*\mathscr{G} \simeq g_0^*\mathscr{G} = \mathscr{H}_0$ which induces a map of sheaves $u : g^*\mathscr{G} \to j_*\mathscr{H}_0 = \mathscr{H}$. Concretely, the map $u$ is given on stalks by the map
\[
(g^*\mathscr{G})_\mu = \mathscr{G}_{g(\mu)} = \int_T G(f_0(t))d(\int_S \nu_s d\mu) \xrightarrow{\ \Delta_{\lambda,\mu_\bullet}\ } \int_S (\int_T G(f_0(t))d\nu_s)d\mu \simeq \mathscr{H}_\mu,
\]
where $\Delta_{\lambda,\mu_\bullet}$ is the categorical Fubini transformation in the category of sets.

Let $h_0 : S \to T$ be the map of sets given by the composition
\[
S \xrightarrow{\ g_0\ } \beta T \xrightarrow{\ f\ } X \xrightarrow{\ f_0^{-1}\ } T \xrightarrow{\ i\ } \beta T
\]
For each $s \in S$, the right ultrastructure on the functor $G$ determines a map
\[
\mathscr{G}_{(i \circ h_0)(s)} \simeq \mathscr{G}_{0h_0(s)} \simeq G((f \circ g_0)(s)) = G(\int_T f_0(t)d\nu_s) \xrightarrow{\ \sigma_{\nu_s}\ } \int_T G(f_0(t))d\nu_s \simeq \mathscr{H}_{0s}
\]
Let $h : \beta S \to \beta T$ denote the unique continuous extension of $h_0$, given concretely by $h(\mu) = h_{0*}(\mu)$. The preceding discussion then gives a map of sheaves $j^*h^*\mathscr{G} \to \mathscr{H}_0$ on the discrete space $S$, which we identify with a map $v : h^*\mathscr{G} \to j_*\mathscr{H}_0 = \mathscr{H}$ on the topological space $\beta S$. Concretely, the map $v$ is given on stalks by
%% [NOTE: Notation 3.5.2 continues across the page break; the sentence "the map v is given on stalks by"
%%  is completed by the display at the top of printed page 40.]

%% --- page 40 ---

the composition
\[
\begin{array}{rcl}
(h^*\mathscr{G})_\mu & \simeq & \int_T G(f_0(t))d(h_{0*}\mu)\\[4pt]
& \xrightarrow{\ \Delta_{\mu,h_0}\ } & \int_S G((f_0 \circ h_0)(s))d\mu\\[4pt]
& = & \int_S G(\int_T f_0(t)d\nu_s)d\mu\\[4pt]
& \xrightarrow{\ \int_S \sigma_{\nu_s}d\mu\ } & \int_S (\int_T G(f_0(t))d\nu_s)d\mu.
\end{array}
\]
\end{notation}

Note that we have a commutative diagram of topological spaces
\[
\beta S \underset{h}{\overset{g}{\rightrightarrows}} \beta T \xrightarrow{\ f\ } X;
\]
that is, we have $f \circ g = f' = f \circ h$ for some continuous map $f' : \beta S \to X$, given concretely by the formula $f'(\mu) = \int_S (\int_T f_0(t)d\nu_s)d\mu$. The maps $u$ and $v$ therefore determine maps $u', v' : f_*(\mathscr{G}) \to f'_*(\mathscr{H})$ in the category $\mathrm{Shv}(X)$.

\renewcommand{\thelemma}{3.5.3}
\begin{lemma}
The sheaf $\mathscr{F}_G$ of Construction 3.4.2 can be identified with the equalizer of the maps $u', v' : f_*(\mathscr{G}) \rightrightarrows f'_*(\mathscr{H})$.
\end{lemma}

\begin{proof}
By construction, the direct image sheaf $f_*\mathscr{G} \simeq f_{0*}\mathscr{G}_0$ is given by
\[
(f_*\mathscr{G})(U) \simeq \prod_{x \in U} G(x).
\]
Let $\mathscr{E}$ denote the equalizer of $u'$ and $v'$, regarded as a subsheaf of $f_*\mathscr{G}$. Unwinding the definitions, we see that for each subset $U \subseteq X$, we can identify $\mathscr{E}(U)$ with the subset of $\prod_{x \in U} G(x)$ consisting of those tuples $\{\phi_x\}_{x \in U}$ which satisfy the following condition:

\smallskip
\noindent $(\ast')$ For every ultrafilter $\nu$ on $T$ such that $\int_T f_0(t)d\nu$ belongs to $U$, the elements $\{\phi_{f_0(t)}\}_{t \in f_0^{-1}(U)}$ and $\phi_{\int_T f_0(t)d\nu}$ have the same image under the maps
\[
\prod_{t \in f_0^{-1}(U)} G(f_0(t)) \xrightarrow{\ q_\nu^{f_0^{-1}(U)}\ } \int_T G(f_0(t))d\nu \xleftarrow{\ \sigma_\nu\ } G(\int_T f_0(t)d\nu),
\]
where $\sigma_\nu$ is given by the left ultrastructure on $G$.
\smallskip

Note that condition $(\ast')$ follows immediately from condition $(\ast)$ of Construction 3.4.2 (applied to the subset $T_0 = f^{-1}(U) \subseteq T$, and the ultrafilter on $T_0$ given by the restriction of $\nu$). That is, we can regard the sheaf $\mathscr{F}_G$ of Construction 3.4.2 as a subsheaf of $\mathscr{E}$. To complete the proof, we must show that $(\ast')$ implies $(\ast)$. To this end, suppose we are given an arbitrary map of sets $e : R \to U$ and an ultrafilter $\lambda$ on $R$ satisfying $\int_R e(r)d\lambda \in U$; we wish to show that $\{\phi_{e(r)}\}_{r \in R}$ and $\phi_{\int_R e(r)d\lambda}$ have the same image under the maps
\[
\prod_{r \in R} G(e(r)) \xrightarrow{\ q_\lambda\ } \int_R G(e(r))d\lambda \xleftarrow{\ \sigma_\lambda\ } G(\int_R e(r)d\lambda).
\]
Since $f_0$ is bijective, the map $e$ factors uniquely as a composition $R \xrightarrow{\ e'\ } T_0 \xrightarrow{\ f_0\ } X$. The desired assertion now follows by applying $(\ast)$ to the ultrafilter $e'_*\lambda$ on $T_0 = f_0^{-1}(U)$.
\end{proof}

Let us now fix a point $y \in X$ and an element $\phi_y$ of the set $G(y)$. For each ultrafilter $\nu$ on $T$ satisfying $f(\nu) = \int_T f_0(t)d\nu = y$, let $\psi_\nu$ denote the image of $\phi_y$ under the map
\[
G(y) = G(\int_T f_0(t)d\nu) \xrightarrow{\ \sigma_\nu\ } \int_T G(f_0(t))d\nu = \mathscr{G}_\nu.
\]

\renewcommand{\thelemma}{3.5.4}
\begin{lemma}
Let $\mu$ be a point of $\beta S$ satisfying $f'(\mu) = y$. Let $u_\mu : \mathscr{G}_{g(\mu)} \to \mathscr{H}_\mu$ and $v_\mu : \mathscr{G}_{h(\mu)} \to \mathscr{H}_\mu$ be the maps induced by $u$ and $v$, respectively. Then $u_\mu(\psi_{g(\mu)}) = v_\mu(\psi_{h(\mu)})$.
\end{lemma}

%% --- page 41 ---

\begin{proof}
This follows from the definition of $u$ and $v$, together with the commutativity of the diagram
%% [NOTE: the source displays the following as one commutative diagram; the array below preserves every
%%  node, arrow and label of that diagram, but its geometric layout is approximated.]
\[
\begin{array}{ccccc}
G(y) & \xrightarrow{\ \sigma_{h_{0*}\mu}\ } & \int_T G(f_0(t))d(h_{0*}\mu) & \xrightarrow{\ \Delta_{\mu,h_0}\ } & \int_S G((f_0 \circ h_0)(s))d\mu\\[10pt]
\Big\Vert & \searrow{\ \ \sigma_\mu} & & & \Big\Vert\\[10pt]
G(\int_T f_0(t)d(\int_S \nu_s d\mu)) & & & & \int_S G(\int_T f_0(t)d\nu_s)d\mu\\[10pt]
& \searrow{\ \ \sigma_{\int_S \nu_s d\mu}} & & & \Big\downarrow{\ \ \int_S \sigma_{\nu_s}d\mu}\\[10pt]
& & \int_T G(f_0(t))d(\int_S \nu_s d\mu) & \xrightarrow{\ \Delta_{\mu,\nu_\bullet}\ } & \int_S (\int_T G(f_0(t))d\nu_s)d\mu.
\end{array}
\]
\end{proof}

\renewcommand{\thelemma}{3.5.5}
\begin{lemma}
There exists a global section $\psi$ of the sheaf $\mathscr{G}|_{f^{-1}\{y\}}$ such that, for each point $\nu \in f^{-1}\{y\}$, the germ of $\psi$ at $\nu$ is equal to $\psi_\nu$.
\end{lemma}

\begin{proof}
Let $\overline{y}$ denote the unique element of $T$ satisfying $f_0(\overline{y}) = y$, and let $K \subseteq \beta S$ denote the inverse image of $\delta_{\overline{y}}$ under the map $h : \beta S \to \beta T$. Then $g : \beta S \to \beta T$ restricts to a map of topological spaces $g_K : K \to f^{-1}\{y\}$. Then $u$ restricts to a map of sheaves $u_K : g_K^*\mathscr{G}|_{f^{-1}\{y\}} \to \mathscr{H}|_K$ on $K$, which we can identify with a map $\iota : \mathscr{G}|_{f^{-1}\{y\}} \to (u_{K*}\mathscr{H}|_K)$ on the fiber $f^{-1}\{y\}$. We first claim that $\iota$ is injective. To prove this, it suffices to observe that for each point $\nu \in f^{-1}\{y\}$, there exists a point $\mu \in K$ for which $u_K(\mu) = \nu$ and $u$ induces a monomorphism of stalks $u_\mu : \mathscr{G}_\mu \to \mathscr{H}_\nu$. In fact, writing $\nu = \nu_s$ for some $s \in S$, we can take $\mu = \delta_s$ to be the principal ultrafilter at the point $s$, in which case the map $u_\mu$ is bijective.

For each point $\nu \in f^{-1}\{y\}$, let $\psi'_\nu$ denote the image of $\psi_\nu$ in the stalk $(u_{K*}\mathscr{H}|_K)_\nu$ (under the map $\iota$). Since $\iota$ is a monomorphism, it will suffice to show that the germs $\{\psi'_\nu\}_{\nu \in f^{-1}\{y\}}$ determine a global section of the sheaf $u_{K*}\mathscr{H}|_K$. Applying Proposition 0.0.10 to the map $u_K : K \to f^{-1}\{y\}$, this can be reformulated as follows:

%% [NOTE: the cross-reference in the preceding sentence is printed as "0.0.10" in the source; it is
%%  transcribed as printed.]
\smallskip
\noindent $(\ast)$ There exists a global section $\xi$ of the sheaf $\mathscr{H}|_K$ having the property that, for every point $\mu \in K$, the germ $\xi_\mu$ of $\xi$ at the point $\mu$ is coincides with the image of $\psi_{g(\mu)}$ under the map $\mathscr{G}_{g(\mu)} \to \mathscr{H}_\mu$ determined by $u_K$.
\smallskip

%% [NOTE: the printed line of $(\ast)$ above reads "is coincides with"; the wording is transcribed as
%%  printed.]
Let $\underline{G(y)}_K$ denote the constant sheaf on $K$ with the value $G(y)$, so that $v$ restricts to a map of sheaves
\[
v_K : \underline{G(y)}_K \to \mathscr{H}|_K.
\]
The map $v_K$ carries the element $\phi_y \in G(y)$ to a global section $\xi$ the sheaf $\mathscr{H}|_K$. It follows from Lemma 3.5.4 that $\xi$ satisfies the requirement of $(\ast)$.
\end{proof}

%% [NOTE: the printed sentence above reads "to a global section \xi the sheaf \mathscr{H}|_K" (without
%%  "of"); it is transcribed as printed.]

\begin{proof}[Proof of Proposition 3.4.6]
Let $\mathscr{G}(f^{-1}\{y\})$ denote the set of global sections of the sheaf $\mathscr{G}|_{f^{-1}\{y\}}$, and define $\mathscr{H}(f'^{-1}\{y\})$ similarly. Since $f$ and $f'$ are proper maps, we can identify $\mathscr{G}(f^{-1}\{y\})$ and $\mathscr{H}(f'^{-1}\{y\})$ with the stalks $(f_*\mathscr{G})_y$ and $(f'_*\mathscr{H})_y$, respectively. Using Lemma 3.5.3 we obtain an equalizer diagram of sets
\[
\mathscr{F}_{G,y} \to \mathscr{G}(f^{-1}\{y\}) \rightrightarrows \mathscr{H}(f'^{-1}\{y\}).
\]
For each element $\phi_y \in G(y)$, the section $\psi \in \mathscr{G}(f^{-1}\{y\})$ of Lemma 3.5.5 belongs to this equalizer (by Lemma 3.5.4), and can therefore be identified with an element of $\mathscr{F}_{G,y}$. It follows immediately from the construction that the canonical map $\theta : \mathscr{F}_{G,y} \to G(y)$ carries this element to $\phi_y$, which shows that $\theta$ is surjective; injectivity was established in Lemma 3.5.1.
\end{proof}

%% [NOTE: the proof of Proposition 3.4.6 ends above. What follows stands on the same supplied page
%%  (printed page 41) and opens section 4; it is transcribed for completeness of the page.]
\smallskip

\noindent\textbf{4. ULTRACATEGORIES AS STACKS}
Let $\mathrm{Comp}$ denote the category of compact Hausdorff spaces. In §3 we showed that every compact Hausdorff space $X$ can be regarded as an ultracategory (Proposition 3.3.1), and that this observation determines a fully faithful embedding from the ordinary category $\mathrm{Comp}$ to the 2-category $\mathbf{Ult}^{\mathrm{L}}$ of Remark 1.4.6 (see Theorem 3.1.5). In this section, we will use embedding as a tool to analyze the entire category $\mathbf{Ult}^{\mathrm{L}}$. To any ultracategory $\mathcal{M}$, we can associate a presheaf of categories on $\mathrm{Comp}$, given by the construction
%% [NOTE: printed page 41 ends here; the sentence continues on printed page 42.]

\end{document}
